\documentclass[12pt,a4paper]{article}
\usepackage[margin=2.5cm]{geometry}
\usepackage{hyperref}
\usepackage{xcolor}
\usepackage{parskip}
\usepackage{booktabs}
\usepackage{enumitem}
\usepackage{array}

\definecolor{futureblue}{RGB}{0,82,204}
\newcommand{\futurework}[1]{\textcolor{futureblue}{\textit{How to solve/integrate: #1}}}

\title{\textbf{Thesis Progress Notes (Code-Verified)}}
\author{Ziv Baretto}
\date{\today}

\begin{document}
\maketitle

\noindent\textbf{Note:} Every statement below is based only on the local code and saved outputs I checked in \path{Fixed_GPU_OpenNyai}, \path{Timeline_Maker}, \path{GRAPH_VISUALISER}, and \path{section_GNN}. Where the repo does not yet implement something, I state that explicitly and add a blue integration note.

\section*{Abstract}

The implemented pipeline converts legal case JSONs into reasoning-focused heterogeneous graphs for outcome prediction. In the data-building stage, \path{Timeline_Maker/merge_cases_v2.py} can group multiple dated JSON files belonging to the same case, deduplicate same-date copies, and produce one merged case document while preserving hearing-level provenance such as \texttt{\_hearing\_timeline}, \texttt{\_section\_offsets}, and sentence-level hearing metadata. In the preprocessing stage, sentence-level OpenNyAI outputs are converted into leakage-safe text sections and normalized legal entities using the fixed-open pipeline in \path{section_GNN/experiments/fixed_open_pipeline/preprocess_fixed_open.py}. The graph builder in \path{section_GNN/final_graph/updated_graph} keeps pre-judgment sections such as \texttt{preamble}, \texttt{facts}, and argument sections, constructs text and entity nodes, and applies temporal gating to statute/provision/precedent citation edges using \texttt{case\_year}. The training code predicts legal outcome from the final \texttt{case} node embedding and evaluates it with saved fold metrics, confusion matrices, and prediction files. Completed outputs include five domain-specific 5-fold runs, a 70{,}460-case cross-bucket 5-fold run, family-domain ablations, cross-bucket text-only and hashing ablations, a fixed-open encoder benchmark, a fixed-open dataset-size sweep, and a separate 3-way outcome experiment. In the saved summaries, the cross-bucket full run reaches accuracy \(0.7854\) and macro-F1 \(0.7785\), the best domain-specific macro-F1 is \(0.7991\) on family/matrimonial, and the fixed-open encoder benchmark shows BGE-M3 as the best average encoder among the tested variants.

\section*{1. What legal document are we processing?}

In the codebase, a case can appear either as a single JSON file or as multiple dated JSON files. The clearest evidence is in \path{DATA_SET_BUILDER_AND_EXPLORER/Timeline_Maker/merge_cases_v2.py}.

\begin{itemize}[noitemsep]
\item Files are expected in the format \texttt{Case Name on DD Month, YYYY.json}.
\item Same-date duplicates are resolved by keeping the file with the most sentences.
\item Multi-date case files are sorted from oldest to latest and merged into one \texttt{\_MERGED.json}.
\item The merged output preserves provenance through \texttt{\_merged\_from}, \texttt{\_hearing\_timeline}, \texttt{\_section\_offsets}, and sentence-level fields \texttt{\_hearing\_date}, \texttt{\_hearing\_file}, and \texttt{\_hearing\_offset}.
\item The latest document supplies the final stored outcome fields: \texttt{case\_outcome\_label}, \texttt{case\_outcome\_score}, and \texttt{llm\_case\_outcome}.
\item Earlier hearings are only appended when their text is not already contained in the latest file.
\end{itemize}

So, in our pipeline, we do \textbf{not} rely on a single guaranteed consolidated document at source. We may start from multiple dated documents for one case and build one merged case document ourselves when needed.

I did \textbf{not} find code or saved documentation here that explains how IndiaKanoon itself constructs or consolidates its document pages. That part is therefore not answered by this repo.

\futurework{If you want this thesis to answer the IndiaKanoon part, add a short separate data-source note based on direct source inspection: collect a few IndiaKanoon examples for the same case across dates, compare page structure, and document whether the site exposes hearing-level snapshots or only a final compiled page.}

\section*{2. Generic case components, graph sections, prediction sections, and timestamps}

\subsection*{What components are visible in the current pipeline?}

The pipeline does not define a manual case-file schema such as ``filed details / evidence / decision'' directly. Instead, it works from sentence-level rhetorical roles. In \path{section_GNN/experiments/fixed_open_pipeline/preprocess_fixed_open.py} and the corresponding README, the important roles are:

\begin{itemize}[noitemsep]
\item \texttt{PREAMBLE}
\item \texttt{FAC}
\item \texttt{ARG\_PETITIONER}
\item \texttt{ARG\_RESPONDENT}
\item \texttt{PRE\_RELIED}, \texttt{PRE\_NOT\_RELIED}, \texttt{STA}
\item \texttt{ANALYSIS}, \texttt{ISSUE}, \texttt{RATIO}, \texttt{RLC}, \texttt{RPC}, \texttt{NONE}
\end{itemize}

From these roles, the current preprocessing builds the following leakage-safe sections:

\begin{itemize}[noitemsep]
\item \texttt{preamble}
\item \texttt{facts}
\item \texttt{arguments}
\item \texttt{petitioner\_arguments}
\item \texttt{respondent\_arguments}
\item \texttt{other\_lawyer\_arguments}
\end{itemize}

There is \textbf{no separate explicit evidence section} in the current section builder. There is also no dedicated prediction target for ``reason given for decision'' in the current code.

\subsection*{Which sections are used for graph building?}

The reasoning-focused graph policy in \path{section_GNN/final_graph/updated_graph/reasoning_graph_policy.py} and the schema in \path{section_GNN/src/graph/schema.py} show that graph construction uses:

\begin{itemize}[noitemsep]
\item Text nodes: \texttt{preamble}, \texttt{facts}, \texttt{arguments}, \texttt{petitioner\_arguments}, \texttt{respondent\_arguments}, \texttt{other\_lawyer\_arguments}
\item Entity nodes: \texttt{petitioner}, \texttt{respondent}, \texttt{court}, \texttt{judge}, \texttt{petitioner\_lawyer}, \texttt{defence\_lawyer}, \texttt{lawyer}, \texttt{statute}, \texttt{provision}, \texttt{precedent}
\end{itemize}

The policy explicitly excludes \texttt{org}, \texttt{gpe}, \texttt{date}, and \texttt{case\_number} from the reasoning-focused graph.

\subsection*{How are those sections extracted from the case documents?}

The extraction is sentence-based:

\begin{itemize}[noitemsep]
\item \texttt{PREAMBLE} sentences, and optionally the text before \texttt{preamble\_end\_char\_offset}, become \texttt{preamble}
\item \texttt{FAC} sentences become \texttt{facts}
\item \texttt{ARG\_PETITIONER} becomes \texttt{petitioner\_arguments} and also contributes to \texttt{arguments}
\item \texttt{ARG\_RESPONDENT} becomes \texttt{respondent\_arguments} and also contributes to \texttt{arguments}
\item \texttt{PRE\_RELIED}, \texttt{PRE\_NOT\_RELIED}, and \texttt{STA} become \texttt{other\_lawyer\_arguments} and also contribute to \texttt{arguments}
\end{itemize}

The following roles are dropped conservatively to avoid leakage into the prediction input: \texttt{ANALYSIS}, \texttt{ISSUE}, \texttt{NONE}, \texttt{RATIO}, \texttt{RLC}, and \texttt{RPC}.

\subsection*{Which sections are used for prediction?}

For the main binary runs, prediction uses \texttt{case\_outcome\_score} as the label field. In the cross-bucket config at \path{section_GNN/runs/cross_bucket_total_dataset/config.yaml}, labels are mapped as:

\begin{itemize}[noitemsep]
\item \texttt{-1} \(\rightarrow\) \texttt{-1}
\item \texttt{0} \(\rightarrow\) \texttt{-1}
\item \texttt{1} \(\rightarrow\) \texttt{1}
\end{itemize}

The classifier predicts from the final \texttt{case} node embedding. The initial text stored on the \texttt{case} node is the concatenation of:

\begin{itemize}[noitemsep]
\item \texttt{preamble}
\item \texttt{facts}
\item \texttt{arguments}
\end{itemize}

This is implemented in \path{section_GNN/final_graph/updated_graph/case_star_builder.py}. The role-specific argument nodes and entity nodes influence the \texttt{case} node through message passing, but the decision/reasoning sections themselves are not given to the model as input.

\subsection*{Are we maintaining timestamps anywhere?}

Yes, but only partially.

\begin{itemize}[noitemsep]
\item In the merged source JSONs, \path{merge_cases_v2.py} preserves full hearing-level provenance and ordering.
\item In model preprocessing, the only temporal value clearly extracted as metadata is \texttt{case\_year}, derived from DATE entities in \path{section_GNN/src/preprocessing/extract.py}.
\item In graph building, \texttt{case\_year} is used for temporal gating of citation edges: precedents must predate the case, and statutes/provisions must not be newer than the case.
\end{itemize}

I did \textbf{not} find any use of full hearing timelines, hearing durations, or sentence-level hearing order as model features in \path{section_GNN}. So the current model is \textbf{not} a time-series model over \(T_1 \ldots T_n\); it uses a single leakage-cleaned case snapshot, plus coarse \texttt{case\_year}.

\futurework{To model true procedural time, carry \texttt{\_hearing\_timeline} from the merged JSONs into cleaned-case metadata, derive features such as hearing count, case duration, and gap statistics, and then add them either as scalar case features or as explicit hearing nodes with temporal edges.}

\section*{3. How are we verifying the neural representations created for the case files?}

Current verification is \textbf{indirect}, through downstream prediction quality rather than a separate representation-analysis stage.

In \path{section_GNN/src/training/train.py}, the training function returns:

\begin{itemize}[noitemsep]
\item split metrics for train/validation/test
\item \texttt{predictions\_df}
\item logits and probabilities
\item \texttt{case\_embeddings}
\end{itemize}

In \path{section_GNN/src/training/metrics.py}, evaluation computes:

\begin{itemize}[noitemsep]
\item accuracy
\item macro-F1 and micro-F1
\item confusion matrix
\item classification report
\item per-class precision, recall, F1, support
\item ROC-AUC and PR-AUC for binary tasks
\end{itemize}

So, as of the checked code and outputs, neural representations are being verified by whether they support accurate legal-outcome classification, not by a dedicated embedding-space study.

I did \textbf{not} find an implemented analysis such as t-SNE/UMAP plots, nearest-neighbour case retrieval, probing classifiers on embeddings, or calibration/cluster checks for the saved \texttt{case\_embeddings}.

\futurework{Use the saved \texttt{case\_embeddings} from \path{train.py} to build a separate representation-analysis notebook or script: export embeddings with labels/domain IDs, run t-SNE or UMAP, inspect nearest neighbours, and train a lightweight probe to test whether domain, label, or citation structure is recoverable from the learned vectors.}

\section*{4. What experiments have we done? What were they trying to test? What are we seeing?}

\subsection*{What the experiment design is testing}

Based on the experiment names, configs, and saved summaries, the completed experiments are testing:

\begin{itemize}[noitemsep]
\item whether performance changes across legal domains
\item whether cross-bucket pooling helps
\item whether cross-case graph sharing helps
\item whether non-text graph structure helps compared with text-only variants
\item whether GNN depth matters
\item whether the text encoder matters when graph/training are held fixed
\item whether increasing training data helps in the fixed-open setup
\item whether the task can be extended from binary to 3-way outcome prediction
\end{itemize}

\subsection*{Completed main 5-fold binary runs}

\begin{center}
\small
\begin{tabular}{lcc}
\toprule
\textbf{Run} & \textbf{Accuracy Mean} & \textbf{Macro-F1 Mean} \\
\midrule
Family/matrimonial & 0.8169 & 0.7991 \\
Financial fraud & 0.7754 & 0.7639 \\
Land/property & 0.8171 & 0.7963 \\
Motor accidents & 0.8319 & 0.7956 \\
Sexual offences & 0.7791 & 0.7526 \\
Cross-bucket total dataset & 0.7854 & 0.7785 \\
\bottomrule
\end{tabular}
\end{center}

For the cross-bucket run, the log at \path{section_GNN/runs/cross_bucket_total_dataset/run_full.log} shows \(73{,}809\) raw files were preprocessed, \(70{,}460\) cleaned cases were written, and \(3{,}349\) files were skipped because labels were missing or unmapped.

\subsection*{Family-domain ablations}

\begin{center}
\small
\begin{tabular}{lcc}
\toprule
\textbf{Configuration} & \textbf{Accuracy Mean} & \textbf{Macro-F1 Mean} \\
\midrule
Full model & 0.8169 & 0.7991 \\
No cross-case sharing & 0.8176 & 0.7966 \\
Text-only & 0.8136 & 0.7944 \\
Depth 1 & 0.8115 & 0.7925 \\
Depth 2 & 0.8173 & 0.7965 \\
Depth 3 & 0.8155 & 0.7978 \\
\bottomrule
\end{tabular}
\end{center}

What we are seeing from these saved outputs:

\begin{itemize}[noitemsep]
\item removing cross-case sharing changes results only slightly in this domain
\item text-only is slightly below the full model
\item one GNN layer is clearly weaker than the 2-layer and 3-layer variants
\item the differences among full, no-cross-case, depth-2, and depth-3 are real but small
\end{itemize}

\subsection*{Cross-bucket ablations}

\begin{center}
\small
\begin{tabular}{lcc}
\toprule
\textbf{Configuration} & \textbf{Accuracy Mean} & \textbf{Macro-F1 Mean} \\
\midrule
Full cross-bucket & 0.7854 & 0.7785 \\
Text-only (variant 1) & 0.7918 & 0.7843 \\
Text-only (variant 2) & 0.7776 & 0.7702 \\
Text-only + hashing & 0.7673 & 0.7594 \\
\bottomrule
\end{tabular}
\end{center}

The important point here is that the saved outputs do \textbf{not} show a single monotonic result in favor of graph structure on the cross-bucket setting. Text-only remains competitive, and one text-only variant is slightly above the full cross-bucket run. What is very clear, however, is that semantic text encoding is better than hashing in this setup.

\subsection*{Controlled encoder benchmark on the fixed-open financial-fraud setup}

\begin{center}
\small
\begin{tabular}{lcc}
\toprule
\textbf{Encoder} & \textbf{Val Macro-F1} & \textbf{Test Macro-F1} \\
\midrule
Hashing & 0.7427 \(\pm\) 0.0067 & 0.7401 \(\pm\) 0.0077 \\
BGE-M3 & 0.7662 \(\pm\) 0.0059 & 0.7513 \(\pm\) 0.0138 \\
InLegalBERT & 0.7412 \(\pm\) 0.0075 & 0.7345 \(\pm\) 0.0016 \\
Legal Longformer & 0.7378 \(\pm\) 0.0037 & 0.7177 \(\pm\) 0.0042 \\
\bottomrule
\end{tabular}
\end{center}

This experiment keeps the graph pipeline fixed and changes only the text encoder. The consistency report in \path{outputs/fin_fraud_labelled_reasoning/multi_embed_test/summary/consistency_checks.json} confirms that split assignments, node counts, relation mappings, and case-split counts are the same across encoder variants. The checked output shows BGE-M3 as the best average encoder among the tested variants.

\subsection*{Other completed experiments with saved outputs}

\begin{itemize}[noitemsep]
\item \textbf{Fixed-open dataset-size sweep:} train sizes \(2500\), \(5000\), and ``all remaining'' were run. Test macro-F1 rises only slightly from \(0.4776\) to \(0.4905\) in the saved summary, so that setup is currently much weaker than the timed-bucket runs.
\item \textbf{3-way outcome experiment:} \path{outputs/augmented_jsons_3way/models/augmented_jsons_3way_reasoning_focused/metrics.json} reports test accuracy \(0.7437\) and macro-F1 \(0.7245\) for labels \(\{-1, 0, 1\}\). This is an outcome-class experiment, not a reason-prediction experiment.
\end{itemize}

\section*{5. Do we have explainability or visibility into the process?}

\subsection*{What currently exists}

Yes, but mostly at the graph-structure and prediction-output level.

\begin{itemize}[noitemsep]
\item \path{GRAPH_VISUALISER/app.py} provides interactive graph views, including a hub network and case-connection exploration.
\item \path{GRAPH_VISUALISER/entity_analysis/analyse.py} computes degree centrality, weighted degree, PageRank, optional betweenness, bridge-entity statistics, and outcome win-rate / mean-score summaries for entities.
\item \path{section_GNN/src/training/train.py} saves \texttt{predictions.csv} with per-case predictions and confidence.
\item \path{section_GNN/final_graph/updated_graph/pipeline.py} stores per-case debug samples including retained text plus node and edge counts.
\item \path{GRAPH_VISUALISER/outputs/stats.json} shows that graph-level descriptive statistics are being produced, such as total nodes, total edges, average degree, maximum degree, and top hubs.
\end{itemize}

\subsection*{What does not yet exist}

I did \textbf{not} find an implemented analysis that answers:

\begin{itemize}[noitemsep]
\item whether the model fails more on long-duration cases
\item whether the model fails more on graphs with many nodes or many citations
\item whether certain courts, buckets, or graph shapes are harder than others
\item which specific edge types or node types receive the most useful model attention during prediction
\end{itemize}

So the honest answer is: \textbf{we have visibility into the constructed graph and final predictions, but we do not yet have a finished failure-analysis layer that explains success/failure by time duration, node count, or similar characteristics.}

\futurework{Join \texttt{predictions.csv} with per-case graph summaries from the reasoning-graph metadata and with hearing metadata from the merged JSONs. Then report error rate by node count, edge count, citation count, hearing count, and case duration. If you want model-level explainability, extend the HGT path to expose attention weights by relation type and summarize them alongside the error analysis.}

\section*{6. Are we predicting only the decision? Can we predict the reasons also?}

For the main thesis pipeline, the answer is \textbf{yes, we are currently predicting only the decision label}. In the main binary runs, that label is derived from \texttt{case\_outcome\_score} and mapped to \(\{-1, 1\}\). There is also a separate saved 3-way setup over \(\{-1, 0, 1\}\), but that is still outcome prediction, not reason prediction.

The current preprocessing explicitly removes direct decision/reasoning content from the model input. In \path{section_GNN/src/graph/schema.py}, forbidden top-level fields include:

\begin{itemize}[noitemsep]
\item \texttt{case\_outcome\_label}
\item \texttt{case\_outcome\_score}
\item \texttt{llm\_case\_outcome}
\item \texttt{decision\_text}
\item \texttt{rpc\_texts}
\end{itemize}

And in the fixed-open preprocessing, the roles \texttt{ANALYSIS}, \texttt{RATIO}, \texttt{RLC}, and \texttt{RPC} are dropped before graph construction. So the current codebase does \textbf{not} train a reason predictor.

I also did \textbf{not} find any implemented co-training or multi-task setup for reasons in the checked folders.

\futurework{If you want to predict reasons, first create reason targets from an explicit taxonomy. One practical route is to normalize \texttt{rpc\_texts} or sentence-level \texttt{RATIO}/\texttt{RPC}/\texttt{ANALYSIS} spans into categories, then train either a multi-label classifier or a second prediction head jointly with the outcome head. Without a reason label schema, the current binary classifier cannot be extended into reason prediction in a controlled way.}

\end{document}
